\documentclass[10pt,a4paper]{amsart}
\usepackage[margin=19mm]{geometry}
\usepackage{amsmath,amssymb,mathtools}
\usepackage[scr=rsfs,cal=euler]{mathalfa}
\usepackage{xfrac}
\usepackage[unicode,hidelinks,bookmarks=false]{hyperref}
\newcommand{\ie}{i.e.\ }
\newcommand{\cf}{cf.\ }
\newcommand{\resp}{respectively\ }
\newcommand{\Subsection}{\subsection}
\providecommand{\nobreakdash}{\nobreak\hbox{-}}
\setlength{\emergencystretch}{2em}
\multlinegap0pt
\usepackage{bm}
\usepackage{bbm}
\usepackage{dsfont}
\usepackage{xspace}
\usepackage{cancel}
\usepackage{mathtools}
\usepackage{amsthm}
\makeatletter
\@ifundefined{theorem}{\newtheorem{theorem}{Theorem}}{}
\makeatother
\title[Regularization of non-overshooting quasi-continuous sliding ]{Regularization of non-overshooting quasi-continuous sliding mode control for chattering suppression at equilibrium}
\author{Michael Ruderman, Denis Efimov}
\date{Source: arXiv:2506.05283v1 (CC BY 4.0)}
\begin{document}
\maketitle

\noindent\emph{Source and licence.} Regularization of non-overshooting quasi-continuous sliding mode control for chattering suppression at equilibrium. Version arXiv:2506.05283v1 (CC BY 4.0), distributed under the Creative Commons Attribution 4.0 International licence, as recorded in the arXiv licence metadata for this exact version and shown on the abstract page https://arxiv.org/abs/2506.05283v1. Full rights notice: licenses/arxiv-2506.05283-CC-BY-4.0.txt.\par\smallskip

\noindent\textit{Editorial scope.} Counted unit: the Theorem whose source label is \texttt{None} in the article (source0.tex lines 528--533, source label \texttt{None}), delivered together with the article's own complete original proof of it (source0.tex lines 534--625, one \texttt{proof} environment of 92 lines). No mathematics is written, repaired or re-derived: statement and proof are the article's own text line for line. Required context retained uncounted: eq:CLoriginal (lines 90--97); eq:33closedloop (lines 229--232); eq:Lyapunov0 (lines 292--294). Every definition, display and label of those lines is retained unchanged. Omitted from this package and not redistributed: 1--89, 98--228, 233--291, 295--527, 626--692. Nothing inside a retained excerpt is omitted or altered: every retained body is delivered exactly as the article's lines are written, so no omission receipt of any kind accompanies this package. Citations: the retained excerpts carry no citation command at all, so no bibliography entry of the article is transcribed and no reference is redistributed. Reference adaptation: none was needed and none was made; every reference inside the delivered unit resolves inside it, so no unresolved reference or citation is produced. Printed slips kept literal: no printed word, symbol, hypothesis or number of the counted statement or of its proof was altered. Layout and class: the article's own class is \texttt{ieeeconf} with packages including graphics, graphicx, epsfig, times, amsmath, amssymb, color, array, lipsum, mathtools, cuted; the wrapper is the lane's pinned \texttt{amsart} editorial wrapper at 10pt on A4, which restates the theorem-like environments the retained text needs and adds the article's own macro definitions that the retained text needs (none), with extra wrapper packages (bm, bbm, dsfont, xspace, cancel, mathtools). The delivered numbering is the unit's own. Assets: no retained span contains a figure environment, a graphics-inclusion command, a PDF-page inclusion, a table or a TikZ picture; no image file is copied into this package, the package declares no asset dependency and no image file is redistributed with it. Licence: version arXiv:2506.05283v1 of this article is distributed under the Creative Commons Attribution 4.0 International licence, as recorded in the arXiv licence metadata for this exact version and shown on the abstract page https://arxiv.org/abs/2506.05283v1. A full rights notice is delivered at licenses/arxiv-2506.05283-CC-BY-4.0.txt. Characters: every retained excerpt is pure ASCII, so the delivered engine, XeLaTeX with its default Latin Modern text font, reports no missing character.
\begin{gather}
\dot{x}_1(t)=x_{2}(t),\label{eq:CLoriginal}\\
\dot{x}_2(t)=\qquad\;\;\nonumber \\
\begin{cases}
-\bigl|x_{1}(t)\bigr|^{-1}\Bigl(\gamma x_{1}(t)+\bigl|x_{2}(t)\bigr|x_{2}(t)\Bigr)+d(t) & x_{1}(t)\ne0\\
-\gamma\bigl|x_{1}(t)\bigr|^{-1}x_{1}(t)+d(t) & x_{1}(t)=0
\end{cases},\nonumber
\end{gather}

\begin{eqnarray}
\dot{x}_{1} & = & x_{2},\label{eq:33closedloop}\\
\dot{x}_{2} & = & -\max\bigl\{\mu,\bigl|x_{1}\bigr|\bigr\}^{-1}\bigl(\gamma x_{1}+\bigl|x_{2}\bigr|x_{2}\bigr)+d,\nonumber 
\end{eqnarray}

\begin{equation}
V(x)=\gamma|x_{1}|+\frac{1}{2}x_{2}^{2}+\varepsilon\sqrt{|x_{1}|}\text{sign}(x_{1})x_{2},\label{eq:Lyapunov0}
\end{equation}

\begin{theorem}
For any choice of $\gamma>\max\{4,2D+4\sqrt{2}D^{1.5}\}$ and $\mu>0$,
with a sufficiently small $D$ (for the chosen value of $\mu$), the system (\ref{eq:33closedloop})
with the disturbances satisfying $\|d\|_{\infty}\leq D$ is strongly
iISS.
\end{theorem}

\begin{proof}
Our strategy to design the required ISS-Lyapunov function consists,
first, in slight modification of one presented in (\ref{eq:Lyapunov0})
by replacing $|x_{1}|$ (integral of nonlinearity in \eqref{eq:CLoriginal})
by $z(x_{1})$ (integral of the counterpart in (\ref{eq:33closedloop})):
\[
\tilde{V}(x)=\gamma z(x_{1})+\frac{1}{2}x_{2}^{2}+\varepsilon\sqrt{z(x_{1})}\text{sign}(x_{1})x_{2},
\]
which is again positive definite and radially unbounded for $\varepsilon\in(0,\sqrt{2\gamma})$.
The time derivative of $\tilde{V}$ on the trajectories of (\ref{eq:33closedloop})
can be written as follows:
\begin{gather*}
\dot{\tilde{V}}=-\frac{|x_{2}|x_{2}^{2}}{\max\{\mu,|x_{1}|\}}-\gamma\varepsilon\frac{\sqrt{z(x_{1})}|x_{1}|}{\max\{\mu,|x_{1}|\}}\\
+\varepsilon\left(\frac{1}{2}-\frac{z(x_{1})}{|x_{1}|}\text{sign}(x_{1})\text{sign}(x_{2})\right)\frac{|x_{1}|}{\sqrt{z(x_{1})}}\frac{x_{2}^{2}}{\max\{\mu,|x_{1}|\}}\\
+(\varepsilon\sqrt{z(x_{1})}\text{sign}(x_{1})+x_{2})d\\
\leq-\frac{|x_{2}|x_{2}^{2}}{\max\{\mu,|x_{1}|\}}-\gamma\varepsilon\frac{\sqrt{z(x_{1})}|x_{1}|}{\max\{\mu,|x_{1}|\}}\\
+\varepsilon\left(\frac{1}{2}+\frac{z(x_{1})}{|x_{1}|}\right)\frac{\frac{|x_{1}|}{\sqrt{z(x_{1})}}x_{2}^{2}}{\max\{\mu,|x_{1}|\}}+(\varepsilon\sqrt{z(x_{1})}+|x_{2}|)|d|.
\end{gather*}
Note that $\frac{z(x_{1})}{|x_{1}|}\leq1$, and
\[
\frac{|x_{1}|}{\sqrt{z(x_{1})}}x_{2}^{2}\leq\frac{2}{3}|x_{2}|^{3}+\frac{1}{3}\left(\frac{|x_{1}|}{\sqrt{z(x_{1})}}\right)^{3}
\]
following the Young's inequality, then we obtain:
\begin{gather*}
\dot{\tilde{V}}\leq(\varepsilon\sqrt{z(x_{1})}+|x_{2}|)|d|-(1-\varepsilon)\frac{|x_{2}|x_{2}^{2}}{\max\{\mu,|x_{1}|\}}\\
-\varepsilon\frac{\left(\gamma\sqrt{z(x_{1})}|x_{1}|-\frac{1}{2}\left(\frac{|x_{1}|}{\sqrt{z(x_{1})}}\right)^{3}\right)}{\max\{\mu,|x_{1}|\}}.
\end{gather*}
Using again the Young's inequality and recalling $|d|\leq D$:
\begin{align*}
|x_{2}||d| & =\frac{|x_{2}|}{\sqrt[3]{\max\{\mu,|x_{1}|\}}}\sqrt[3]{\max\{\mu,|x_{1}|\}}|d|\\
 & \leq\frac{1}{3}\frac{|x_{2}|^{3}}{\max\{\mu,|x_{1}|\}}+\frac{2}{3}\sqrt{\max\{\mu,|x_{1}|\}}D^{1.5},
\end{align*}
the upper estimate for the derivative of the Lyapunov function can
be further simplified:
\begin{gather*}
\dot{\tilde{V}}\leq-\varepsilon k(x_{1})-\left(\frac{2}{3}-\varepsilon\right)\frac{|x_{2}|x_{2}^{2}}{\max\{\mu,|x_{1}|\}} -\varepsilon e(x_{1},D),
\end{gather*}
where
\begin{gather*}
k(x_{1})=\frac{\gamma\sqrt{z(x_{1})}|x_{1}|-\left(\frac{|x_{1}|}{\sqrt{z(x_{1})}}\right)^{3}}{2\max\{\mu,|x_{1}|\}}\\
=\begin{cases}
\gamma\frac{1}{\left(\sqrt{2\mu}\right)^{3}}x_{1}^{2}-\sqrt{2\mu} & |x_{1}|<\mu\\
\frac{\gamma\sqrt{|x_{1}|-\frac{\mu}{2}}|x_{1}|-\left(\frac{|x_{1}|}{\sqrt{|x_{1}|-\frac{\mu}{2}}}\right)^{3}}{2|x_{1}|} & |x_{1}|\geq\mu
\end{cases},\\
e(s,D)=\frac{\sqrt{z(s)}}{\max\{\mu,|s|\}}\left(\frac{\gamma}{2}|s|-D-\frac{\max\{\mu,|s|\}^{\frac{3}{2}}}{\frac{3}{2}\varepsilon\sqrt{z(s)}}D^{\frac{3}{2}}\right).
\end{gather*}
First, analyzing the expression of $k(x_{1})$ we conclude that if
$\gamma>4$, then $k(x_{1})>0$ for $|x_{1}|\geq\mu$ and, while for
$|x_{1}|<\mu$ the constant negative term is obviously dominating
for $|x_{1}|\leq\frac{2\mu}{\sqrt{\gamma}}$, and increasing the gain
$\gamma$ this domain can be narrowed. Performing similar analysis
we get for $|x_{1}|\geq\mu$:
\begin{align*}
e(x_{1},D) & =\sqrt{|x_{1}|-\frac{\mu}{2}}\left(\frac{\gamma}{2}-D-\frac{2\sqrt{|x_{1}|}}{3\varepsilon\sqrt{|x_{1}|-\frac{\mu}{2}}}D^{1.5}\right)\\
 & \geq\sqrt{|x_{1}|-\frac{\mu}{2}}\left(\frac{\gamma}{2}-D-\frac{2\sqrt{2}}{3\varepsilon}D^{1.5}\right)
\end{align*}
that is strictly positive provided that 
\[
\frac{\gamma}{2}>D+\frac{2\sqrt{2}}{3\varepsilon}D^{1.5},
\]
and for $|x_{1}|<\mu$:
\begin{align*}
e(x_{1},D) & =\frac{\frac{\gamma}{2}|x_{1}|}{\mu}\sqrt{\frac{x_{1}^{2}}{2\mu}}-\sqrt{\frac{x_{1}^{2}}{2\mu}}D-\frac{2\sqrt{\mu}}{3\varepsilon}D^{1.5}\\
 & \geq-\sqrt{\frac{\mu}{2}}D-\frac{2\sqrt{\mu}}{3\varepsilon}D^{1.5}
\end{align*}
having the disturbance gain of order $\sqrt{\mu}$. Note that for
$\gamma>4$, $\varepsilon=\frac{1}{3}$ is an admissible choice, giving
the restriction stated in the formulation of the theorem, then $\dot{\tilde{V}}$
is strictly negative for $|x_{1}|\geq\mu$, and locally close to the
origin, with $|x_{1}|<\mu$, a bias and influence of the disturbances
appear. Therefore, $\tilde{V}$ is a practical ISS-Lyapunov function
for (\ref{eq:33closedloop}) with $\|d\|_{\infty}\leq D$. In order
to construct an ISS-Lyapunov function, consider the following candidate:
\[
U(x)=\tilde{V}^{3}(x)+\sigma(W(x)),
\]
where $\sigma\in\mathcal{K}$ is a bounded continuously differentiable
function which is reduced to a linear map close to the origin (it is designed in a way 
to guarantee that $\sigma(W(x)))=\text{const}$ for $|x_1|\geq\mu$). The
idea behind this design is as follows. First, clearly, $U$ is positive
definite and radially unbounded due to $\tilde{V}$ possesses these
properties. Second, the derivative of the term $\tilde{V}^{3}(x)$
has the form $3\tilde{V}^{2}\dot{\tilde{V}}$, and it is strictly
negative for $|x_{1}|\geq\mu$, and where $\dot{U}=3\tilde{V}^{2}\dot{\tilde{V}}$
due to a choice of $\sigma$, while for $|x_{1}|<\mu$ the bias is
now multiplied by the term of orders $x_{1}^{4}$ and $x_{2}^{4}$
which can be compensated by the negative terms of order $x_{1}^{4}$
and $|x_{2}|x_{2}^{2}$ contained in $\dot{W}$, provided that the
disturbances are sufficiently small. Hence, by a proper weighting
of $\tilde{V}$ and $W$, we can guarantee that $U$ is an ISS-Lyapunov
function for (\ref{eq:33closedloop}) with $\|d\|_{\infty}\leq D$.
\end{proof}

\end{document}
